% 2.4 晶体学点群的对称适配函数
\section{晶体学点群的对称适配函数}

在把所有点群的结果列成表之前，还要处理一个理论问题：由式 (2.3.2) 对给定表示的给定行（即固定 $l$ 与 $t$）生成的面谐函数不一定相互正交。出于实用考虑，同一表示的任何两组基最好由相互正交的函数组成。下面各表给出的所有展开都已借助定理 2.2.2 与 2.2.3 正交化。

表 2.2 给出 32 个晶体学点群（单值）表示的特征标表。表示按 Mulliken（1933）的记号标记，本书即沿用这一记号；同时并列给出 Koster, Dimmock, Wheeler, and Statz（1963）等使用的 $\Gamma$ 记号以备参照。表 2.3 给出我们对点群简并表示所用的矩阵。2.3 节的方法可用于确定循环群、二面体群与立方群的面谐函数，这些函数分别列于表 2.4--2.6（Altmann 1957，Altmann and Bradley 1963\textsuperscript{b}，Altmann and Cracknell 1965）。下面给出若干解读这些面谐函数表的例子。

\begin{figure}[H]
\centering
\includegraphics[width=0.86\textwidth]{c2_t22_p1.png}
\caption*{表 2.2（第一部分）：$1\ (C_{1})$、$\bar{1}\ (C_{i})$、$2\ (C_{2})$、$m\ (C_{1h})$、$mm2\ (C_{2v})$、$222\ (D_{2})$、$2/m\ (C_{2h})$ 的特征标表（原书扫描占位）。}
\end{figure}
\begin{figure}[H]
\centering
\includegraphics[width=0.86\textwidth]{c2_t22_p2.png}
\caption*{表 2.2（第二部分）：$4\ (C_{4})$、$\bar{4}\ (S_{4})$、$4/m\ (C_{4h})$、$3\ (C_{3})$、$\bar{3}\ (C_{3i})$、$32\ (D_{3})$、$3m\ (C_{3v})$、$\bar{3}m\ (D_{3d})$、$6\ (C_{6})$、$\bar{6}\ (C_{3h})$、$6/m\ (C_{6h})$ 的特征标表（原书扫描占位）。}
\end{figure}
\begin{figure}[H]
\centering
\includegraphics[width=0.86\textwidth]{c2_t22_p3.png}
\caption*{表 2.2（第三部分）：$422\ (D_{4})$、$4mm\ (C_{4v})$、$\bar{4}2m\ (D_{2d})$、$4/mmm\ (D_{4h})$、$622\ (D_{6})$、$6mm\ (C_{6v})$、$\bar{6}2m\ (D_{3h})$、$6/mmm\ (D_{6h})$、$23\ (T)$、$m3\ (T_{h})$、$432\ (O)$、$\bar{4}3m\ (T_{d})$、$m3m\ (O_{h})$ 的特征标表，以及表 2.3（简并表示的矩阵）的开头（原书扫描占位）。}
\end{figure}
\begin{figure}[H]
\centering
\includegraphics[width=0.86\textwidth]{c2_t23.png}
\caption*{表 2.3（续）：四方群、三方与六方群简并表示的矩阵（原书扫描占位）。}
\end{figure}
\begin{figure}[H]
\centering
\includegraphics[width=0.86\textwidth]{c2_t24_t25.png}
\caption*{表 2.3（续）：三方与六方群的矩阵 Key，以及立方群二重简并表示 $E$ 的矩阵（原书扫描占位）。}
\end{figure}
\begin{figure}[H]
\centering
\includegraphics[width=0.86\textwidth]{c2_t25_t26.png}
\caption*{表 2.3（结尾）与表 2.4（循环群的面谐函数）开头（原书扫描占位）。}
\end{figure}
\begin{figure}[H]
\centering
\includegraphics[width=0.86\textwidth]{c2_t26.png}
\caption*{表 2.4（续）：循环群的面谐函数（原书扫描占位）。}
\end{figure}
\begin{figure}[H]
\centering
\includegraphics[width=0.86\textwidth]{c2_t27_p1.png}
\caption*{表 2.4（结尾）及其注，与表 2.5（二面体群的面谐函数）（原书扫描占位）。}
\end{figure}
\begin{figure}[H]
\centering
\includegraphics[width=0.86\textwidth]{c2_t27_p2.png}
\caption*{表 2.5（续）与表 2.6 的开头：三方、六方及立方群的面谐函数（原书扫描占位）。}
\end{figure}
\begin{figure}[H]
\centering
\includegraphics[width=0.86\textwidth]{c2_t27_p3.png}
\caption*{表 2.6（续）：立方群与三方、六方群的面谐函数（原书扫描占位）。}
\end{figure}
\begin{figure}[H]
\centering
\includegraphics[width=0.86\textwidth]{c2_t27_p4.png}
\caption*{表 2.7：$O(3)$ 的表示 $\mathscr{D}^{l}$ 与点群表示的相容表（第一部分，原书扫描占位）。}
\end{figure}
\begin{figure}[H]
\centering
\includegraphics[width=0.86\textwidth]{c2_t27_p5.png}
\caption*{表 2.7（第二部分，原书扫描占位）。}
\end{figure}
\begin{figure}[H]
\centering
\includegraphics[width=0.86\textwidth]{c2_t27_p6.png}
\caption*{表 2.7（第三部分，原书扫描占位）。}
\end{figure}


\textit{例} 2.4.3. 取 $l = 2$（即五重简并的原子 D 项），点群 $mm2$（$C_{2v}$）为 $\mathbf{G}$。按表 2.7，该项在 $mm2$（$C_{2v}$）中劈裂为五个非简并能级：

\begin{equation*}
\mathscr{D}^{2}\{R(\alpha, \beta, \gamma)\} = 2A_{1} + A_{2} + B_{1} + B_{2}.
\tag{2.4.4}
\end{equation*}

$\mathscr{D}^{2}\{R(\alpha, \beta, \gamma)\}$ 的基函数是 $\{Y_{2}^{-2}(\theta, \phi), Y_{2}^{-1}(\theta, \phi), Y_{2}^{0}(\theta, \phi), Y_{2}^{1}(\theta, \phi), Y_{2}^{2}(\theta, \phi)\}$；由表 2.5 可见，这些函数按式 (2.4.4) 分配到 $mm2$（$C_{2v}$）的五个表示：$A_{1}$：$Y_{2}^{0}(\theta, \phi)$，$Y_{2}^{2,c}(\theta, \phi)$；$A_{2}$：$Y_{2}^{2,s}(\theta, \phi)$；$B_{1}$：$Y_{2}^{1,c}(\theta, \phi)$；$B_{2}$：$Y_{2}^{1,s}(\theta, \phi)$。见图 2.2。

\begin{figure}[H]
\centering
\includegraphics[width=0.46\textwidth]{fig2_2_mu.jpg}
\caption{原子 D 项在具有 $mm2$（$C_{2v}$）对称性的晶体场中的示意劈裂。}
\end{figure}

\textit{例} 2.4.4. 作为第二个例子，取 $l = 3$（即七重简并的原子 F 项），点群取立方群 432（$O$）。按表 2.7 有

\begin{equation*}
\mathscr{D}^{3}\{R(\alpha, \beta, \gamma)\} = A_{2} + T_{1} + T_{2}.
\tag{2.4.5}
\end{equation*}

$A_{2}$ 是 432（$O$）的非简并表示，$T_{1}$ 与 $T_{2}$ 是 432（$O$）的三重简并表示。与这些能级对应的波函数角部分是 $Y_{3}^{0}(\theta, \phi)$，$Y_{3}^{1,c}(\theta, \phi)$，$Y_{3}^{1,s}(\theta, \phi)$，$Y_{3}^{2,c}(\theta, \phi)$，$Y_{3}^{2,s}(\theta, \phi)$，$Y_{3}^{3,c}(\theta, \phi)$，$Y_{3}^{3,s}(\theta, \phi)$ 的线性组合。由表 2.6：属于 $A_{2}$ 的面谐函数为 $Y_{3}^{2,s}(\theta, \phi)$；属于 $T_{1}$ 的为 $[\{aY_{3}^{1,c}(\theta, \phi) - bY_{3}^{3,c}(\theta, \phi)\}, \{aY_{3}^{1,s}(\theta, \phi) + bY_{3}^{3,s}(\theta, \phi)\}, \{-Y_{3}^{0}(\theta, \phi)\}]$；属于 $T_{2}$ 的为 $[\{-bY_{3}^{1,c}(\theta, \phi) - aY_{3}^{3,c}(\theta, \phi)\}, \{bY_{3}^{1,s}(\theta, \phi) - aY_{3}^{3,s}(\theta, \phi)\}, \{Y_{3}^{2,c}(\theta, \phi)\}]$，其中 $a = 0.612\,372\,435\,7$，$b = 0.790\,569\,415\,0$。原子 F 项在 432（$O$）晶体场中的劈裂见图 2.3。

\begin{figure}[H]
\centering
\includegraphics[width=0.46\textwidth]{fig2_3_mu.jpg}
\caption{原子 F 项在具有 432（$O$）对称性的晶体场中的示意劈裂（括号内给出各能级的简并度）。}
\end{figure}
